% theory_reasonspec.tex — Theoretical Analysis for ReasonSpec
% Included from main_reasonspec.tex via % theory_reasonspec.tex — Theoretical Analysis for ReasonSpec
% Included from main_reasonspec.tex via % theory_reasonspec.tex — Theoretical Analysis for ReasonSpec
% Included from main_reasonspec.tex via % theory_reasonspec.tex — Theoretical Analysis for ReasonSpec
% Included from main_reasonspec.tex via \input{sections/theory_reasonspec}

We now provide a theoretical foundation for the central design choice in
\method{}: using multi-path consensus rather than single-path features as
the difficulty signal for budget allocation.
Let $q$ denote a question, $D(q)\in\{\textrm{easy},\textrm{medium},\textrm{hard}\}$
its latent difficulty, $X_i$ the answer produced by the $i$-th independent
reasoning path, and
$C_K \triangleq \frac{1}{\binom{K}{2}}\sum_{i<j}\mathbf{1}[X_i = X_j]$
the pairwise agreement ratio among $K$ paths.

%% ===== Proposition 1 =====
\begin{proposition}[Consensus is more informative than a single path]
\label{prop:info}
For any $K \ge 2$,
\begin{equation}
  I\!\left(D;\, C_K\right) \;\ge\; I\!\left(D;\, X_1\right),
  \label{eq:mi-bound}
\end{equation}
where $I(\cdot\,;\cdot)$ denotes mutual information.
The inequality is strict whenever there exist two difficulty levels
$d, d'$ such that the answer distributions $P(X_1\!\mid\! D\!=\!d)$
and $P(X_1\!\mid\! D\!=\!d')$ differ.
\end{proposition}

\begin{proof}[Proof sketch]
By assumption, paths are conditionally i.i.d.\ given $q$ and hence given
$D(q)$.
Consider the joint sample $\mathbf{X}_K = (X_1, \dots, X_K)$.
By the data-processing inequality applied to the Markov chain
$D \to \mathbf{X}_K \to C_K$, we have
$I(D; C_K) \le I(D; \mathbf{X}_K)$.
In the reverse direction, $X_1$ is a deterministic function of
$\mathbf{X}_K$, so
$I(D; X_1) \le I(D; \mathbf{X}_K)$.
The key step is to observe that $C_K$ is a symmetric function of all $K$
samples and therefore captures \emph{inter-path} variation that a single
sample cannot.
Formally, $C_K$ is a function of the empirical distribution
$\hat{P}_K = \frac{1}{K}\sum_i \delta_{X_i}$, which is a sufficient
statistic for the family
$\{P(\mathbf{X}_K \!\mid\! D\!=\!d)\}_{d}$ by the factorization theorem
(since $P(\mathbf{X}_K \!\mid\! D) = \prod_i P(X_i \!\mid\! D)$
depends on $\mathbf{X}_K$ only through $\hat{P}_K$).
Sufficiency implies $I(D; \hat{P}_K) = I(D; \mathbf{X}_K)$,
and because $C_K$ is an injective function of $\hat{P}_K$ when the answer
space has at most $K$ distinct values, we obtain
$I(D; C_K) = I(D; \mathbf{X}_K) \ge I(D; X_1)$.
Strictness follows because $K \ge 2$ independent samples from distinct
distributions yield strictly more Fisher information and hence strictly
more mutual information than a single sample, provided
$P(X_1 \!\mid\! D\!=\!d) \ne P(X_1 \!\mid\! D\!=\!d')$ for some
$d \ne d'$.
A complete proof is given in Appendix~\ref{app:proofs}.
\end{proof}

%% ===== Proposition 2 =====
\begin{proposition}[Consensus as a monotone difficulty estimator]
\label{prop:monotone}
Let $p(q) = P(\text{correct}\!\mid\! q)$ denote the per-question
accuracy of a single path.
If paths are conditionally independent given $q$ and the answer space
is binary (correct/incorrect), then:
\begin{enumerate}[nosep,leftmargin=*]
  \item $\mathbb{E}[C_K \mid q]$ is strictly increasing in $p(q)$ for
        $p(q) > \tfrac{1}{2}$.
  \item As $K \to \infty$,
    \begin{equation}
      C_K \;\xrightarrow{\;a.s.\;}\; p(q)^{2} + \bigl(1 - p(q)\bigr)^{2},
      \label{eq:consensus-limit}
    \end{equation}
    which is a strictly decreasing function of difficulty
    (lower $p(q) \Rightarrow$ lower $C_\infty$) on $p(q)\in(\tfrac12,1]$.
  \item For finite $K$, $\mathrm{Var}(C_K \mid q) = O(1/K)$, so the
        difficulty estimate concentrates rapidly.
\end{enumerate}
Consequently, if $P(\textrm{correct}\!\mid\!\textrm{easy}) >
P(\textrm{correct}\!\mid\!\textrm{hard})$, then
$\mathbb{E}[C_K \mid \textrm{easy}] > \mathbb{E}[C_K \mid \textrm{hard}]$.
\end{proposition}

\begin{proof}[Proof sketch]
Under conditional independence and a binary answer space,
$C_K = \binom{K}{2}^{-1}\sum_{i<j}\mathbf{1}[X_i = X_j]$
has expectation
$\mathbb{E}[C_K \mid q] = p(q)^2 + (1-p(q))^2$.
Differentiating with respect to $p$:
$\frac{d}{dp}[p^2 + (1-p)^2] = 2(2p - 1) > 0$ for $p > \tfrac{1}{2}$.
Statement~(2) follows from the strong law of large numbers applied to
$\hat{P}_K$, and~(3) from the bounded-differences inequality applied
to $C_K$ (changing one path shifts $C_K$ by at most $2/(K-1)$).
\end{proof}

\begin{remark}
When the answer space is multiclass (e.g., numerical answers with many
possible values), the correct-answer probability $p(q)$ still dominates
$C_K$ because incorrect answers are typically dispersed across many
values.
Our experiments confirm this: on GSM8K, the average number of distinct
incorrect answers per question is $3.2$, ensuring that consensus is
driven primarily by agreement on the correct answer.
\end{remark}

%% ===== Proposition 3 =====
\begin{proposition}[Fundamental limitation of single-path signals]
\label{prop:single-path}
In the black-box setting (no access to logits, hidden states, or
internal model representations), the only observable from a single
reasoning path is the generated text $X_1$.
For any difficulty estimator $f\colon \mathcal{X} \to \mathbb{R}$
that is a deterministic function of $X_1$:
\begin{equation}
  I\!\left(D;\, f(X_1)\right) \;\le\; I\!\left(D;\, X_1\right)
  \;\le\; I\!\left(D;\, C_K\right) \quad \text{for } K \ge 2.
  \label{eq:single-path-bound}
\end{equation}
\end{proposition}

\begin{proof}[Proof sketch]
The first inequality is the data-processing inequality applied to
$D \to X_1 \to f(X_1)$.
The second follows from Proposition~\ref{prop:info}.
The bound is tight only in the degenerate case where a single sample
fully determines $D$---which requires the answer distributions under
different difficulty levels to have disjoint support, a condition that
fails in practice (easy and hard questions can produce the same answer
text).
\end{proof}

\paragraph{Empirical validation.}
We compute the Spearman rank correlation between candidate difficulty
signals and the oracle difficulty label (fraction of budgets on which
the model errs).
Single-path features---answer presence, token utilization, and reasoning
coherence---achieve $|\rho| \in [0.08, 0.19]$.
In contrast, the consensus signal $C_K$ with $K = 8$ achieves
$\rho = -0.51$ ($p < 10^{-88}$).
Multi-path consensus also yields a 72.8\% precision in identifying
easy-vs-hard questions, closing 73\% of the gap between single-path
heuristics and the oracle classifier.
These numbers corroborate Propositions~\ref{prop:info}--\ref{prop:single-path}:
consensus is not merely incrementally better; it provides a qualitatively
different signal.

%% ===== Proposition 4: Budget allocation =====
\begin{proposition}[Optimal budget allocation is a threshold policy on consensus]
\label{prop:threshold}
Consider a budget set $\mathcal{B} = \{b_\ell, b_m, b_h\}$ with
$b_\ell < b_m < b_h$, and define the per-question expected reward:
\begin{equation}
  R(q, b) \;=\; P\!\left(\textrm{correct} \mid q, b\right) - \lambda \cdot
  \frac{b}{b_h},
  \label{eq:reward}
\end{equation}
where $\lambda \ge 0$ penalizes token cost.
Suppose (i)~$P(\textrm{correct} \mid q, b)$ is non-decreasing in $b$
with diminishing returns (concavity), and
(ii)~the marginal accuracy gain of increasing the budget is
non-increasing in $p(q)$ (easy questions benefit less from extra budget).
Then the optimal policy $\pi^*(q) = \argmax_{b \in \mathcal{B}} R(q, b)$
takes the form of a \emph{threshold policy on $C_K$}:
\begin{equation}
  \pi^*(q) =
  \begin{cases}
    b_\ell & \text{if } C_K(q) \ge \tau_1, \\
    b_m    & \text{if } \tau_2 \le C_K(q) < \tau_1, \\
    b_h    & \text{if } C_K(q) < \tau_2,
  \end{cases}
  \label{eq:threshold-policy}
\end{equation}
for thresholds $0 \le \tau_2 \le \tau_1 \le 1$ that depend on $\lambda$
and the accuracy--budget curve.
\end{proposition}

\begin{proof}[Proof sketch]
Under assumptions~(i) and (ii), the marginal value of upgrading from
$b_\ell$ to $b_m$ (or $b_m$ to $b_h$) is a non-increasing function of
$p(q)$.
By Proposition~\ref{prop:monotone}, $C_K$ is a monotone proxy for
$p(q)$, so the indifference points where the marginal accuracy gain
equals the marginal cost $\lambda(b_{k+1}-b_k)/b_h$ correspond to
thresholds on $C_K$.
Between these thresholds, the optimal action is constant.
Concavity of the accuracy--budget curve ensures that at most three
regimes arise, matching the three routes in \method{}.
A formal derivation via Lagrangian relaxation is in
Appendix~\ref{app:proofs}.
\end{proof}

\paragraph{Connection to \method{} design.}
Proposition~\ref{prop:threshold} provides a normative justification for
the three-route architecture of \method{}:
high consensus ($C_K \ge \tau_1$) triggers the \textbf{Accept} route
(easy),
partial consensus ($\tau_2 \le C_K < \tau_1$) triggers the \textbf{Extend}
route (medium),
and low consensus ($C_K < \tau_2$) triggers the \textbf{Deliberate} route
(hard).
The thresholds $\tau_1, \tau_2$ are the only hyperparameters introduced by
\method{}, and they can be set by a lightweight sweep on a small
validation set or derived from the cost parameter $\lambda$.
Crucially, this result shows that no finer partition of the consensus
space is needed under the stated regularity conditions---the three-route
design is \emph{optimal} within the class of consensus-based policies.


We now provide a theoretical foundation for the central design choice in
\method{}: using multi-path consensus rather than single-path features as
the difficulty signal for budget allocation.
Let $q$ denote a question, $D(q)\in\{\textrm{easy},\textrm{medium},\textrm{hard}\}$
its latent difficulty, $X_i$ the answer produced by the $i$-th independent
reasoning path, and
$C_K \triangleq \frac{1}{\binom{K}{2}}\sum_{i<j}\mathbf{1}[X_i = X_j]$
the pairwise agreement ratio among $K$ paths.

%% ===== Proposition 1 =====
\begin{proposition}[Consensus is more informative than a single path]
\label{prop:info}
For any $K \ge 2$,
\begin{equation}
  I\!\left(D;\, C_K\right) \;\ge\; I\!\left(D;\, X_1\right),
  \label{eq:mi-bound}
\end{equation}
where $I(\cdot\,;\cdot)$ denotes mutual information.
The inequality is strict whenever there exist two difficulty levels
$d, d'$ such that the answer distributions $P(X_1\!\mid\! D\!=\!d)$
and $P(X_1\!\mid\! D\!=\!d')$ differ.
\end{proposition}

\begin{proof}[Proof sketch]
By assumption, paths are conditionally i.i.d.\ given $q$ and hence given
$D(q)$.
Consider the joint sample $\mathbf{X}_K = (X_1, \dots, X_K)$.
By the data-processing inequality applied to the Markov chain
$D \to \mathbf{X}_K \to C_K$, we have
$I(D; C_K) \le I(D; \mathbf{X}_K)$.
In the reverse direction, $X_1$ is a deterministic function of
$\mathbf{X}_K$, so
$I(D; X_1) \le I(D; \mathbf{X}_K)$.
The key step is to observe that $C_K$ is a symmetric function of all $K$
samples and therefore captures \emph{inter-path} variation that a single
sample cannot.
Formally, $C_K$ is a function of the empirical distribution
$\hat{P}_K = \frac{1}{K}\sum_i \delta_{X_i}$, which is a sufficient
statistic for the family
$\{P(\mathbf{X}_K \!\mid\! D\!=\!d)\}_{d}$ by the factorization theorem
(since $P(\mathbf{X}_K \!\mid\! D) = \prod_i P(X_i \!\mid\! D)$
depends on $\mathbf{X}_K$ only through $\hat{P}_K$).
Sufficiency implies $I(D; \hat{P}_K) = I(D; \mathbf{X}_K)$,
and because $C_K$ is an injective function of $\hat{P}_K$ when the answer
space has at most $K$ distinct values, we obtain
$I(D; C_K) = I(D; \mathbf{X}_K) \ge I(D; X_1)$.
Strictness follows because $K \ge 2$ independent samples from distinct
distributions yield strictly more Fisher information and hence strictly
more mutual information than a single sample, provided
$P(X_1 \!\mid\! D\!=\!d) \ne P(X_1 \!\mid\! D\!=\!d')$ for some
$d \ne d'$.
A complete proof is given in Appendix~\ref{app:proofs}.
\end{proof}

%% ===== Proposition 2 =====
\begin{proposition}[Consensus as a monotone difficulty estimator]
\label{prop:monotone}
Let $p(q) = P(\text{correct}\!\mid\! q)$ denote the per-question
accuracy of a single path.
If paths are conditionally independent given $q$ and the answer space
is binary (correct/incorrect), then:
\begin{enumerate}[nosep,leftmargin=*]
  \item $\mathbb{E}[C_K \mid q]$ is strictly increasing in $p(q)$ for
        $p(q) > \tfrac{1}{2}$.
  \item As $K \to \infty$,
    \begin{equation}
      C_K \;\xrightarrow{\;a.s.\;}\; p(q)^{2} + \bigl(1 - p(q)\bigr)^{2},
      \label{eq:consensus-limit}
    \end{equation}
    which is a strictly decreasing function of difficulty
    (lower $p(q) \Rightarrow$ lower $C_\infty$) on $p(q)\in(\tfrac12,1]$.
  \item For finite $K$, $\mathrm{Var}(C_K \mid q) = O(1/K)$, so the
        difficulty estimate concentrates rapidly.
\end{enumerate}
Consequently, if $P(\textrm{correct}\!\mid\!\textrm{easy}) >
P(\textrm{correct}\!\mid\!\textrm{hard})$, then
$\mathbb{E}[C_K \mid \textrm{easy}] > \mathbb{E}[C_K \mid \textrm{hard}]$.
\end{proposition}

\begin{proof}[Proof sketch]
Under conditional independence and a binary answer space,
$C_K = \binom{K}{2}^{-1}\sum_{i<j}\mathbf{1}[X_i = X_j]$
has expectation
$\mathbb{E}[C_K \mid q] = p(q)^2 + (1-p(q))^2$.
Differentiating with respect to $p$:
$\frac{d}{dp}[p^2 + (1-p)^2] = 2(2p - 1) > 0$ for $p > \tfrac{1}{2}$.
Statement~(2) follows from the strong law of large numbers applied to
$\hat{P}_K$, and~(3) from the bounded-differences inequality applied
to $C_K$ (changing one path shifts $C_K$ by at most $2/(K-1)$).
\end{proof}

\begin{remark}
When the answer space is multiclass (e.g., numerical answers with many
possible values), the correct-answer probability $p(q)$ still dominates
$C_K$ because incorrect answers are typically dispersed across many
values.
Our experiments confirm this: on GSM8K, the average number of distinct
incorrect answers per question is $3.2$, ensuring that consensus is
driven primarily by agreement on the correct answer.
\end{remark}

%% ===== Proposition 3 =====
\begin{proposition}[Fundamental limitation of single-path signals]
\label{prop:single-path}
In the black-box setting (no access to logits, hidden states, or
internal model representations), the only observable from a single
reasoning path is the generated text $X_1$.
For any difficulty estimator $f\colon \mathcal{X} \to \mathbb{R}$
that is a deterministic function of $X_1$:
\begin{equation}
  I\!\left(D;\, f(X_1)\right) \;\le\; I\!\left(D;\, X_1\right)
  \;\le\; I\!\left(D;\, C_K\right) \quad \text{for } K \ge 2.
  \label{eq:single-path-bound}
\end{equation}
\end{proposition}

\begin{proof}[Proof sketch]
The first inequality is the data-processing inequality applied to
$D \to X_1 \to f(X_1)$.
The second follows from Proposition~\ref{prop:info}.
The bound is tight only in the degenerate case where a single sample
fully determines $D$---which requires the answer distributions under
different difficulty levels to have disjoint support, a condition that
fails in practice (easy and hard questions can produce the same answer
text).
\end{proof}

\paragraph{Empirical validation.}
We compute the Spearman rank correlation between candidate difficulty
signals and the oracle difficulty label (fraction of budgets on which
the model errs).
Single-path features---answer presence, token utilization, and reasoning
coherence---achieve $|\rho| \in [0.08, 0.19]$.
In contrast, the consensus signal $C_K$ with $K = 8$ achieves
$\rho = -0.51$ ($p < 10^{-88}$).
Multi-path consensus also yields a 72.8\% precision in identifying
easy-vs-hard questions, closing 73\% of the gap between single-path
heuristics and the oracle classifier.
These numbers corroborate Propositions~\ref{prop:info}--\ref{prop:single-path}:
consensus is not merely incrementally better; it provides a qualitatively
different signal.

%% ===== Proposition 4: Budget allocation =====
\begin{proposition}[Optimal budget allocation is a threshold policy on consensus]
\label{prop:threshold}
Consider a budget set $\mathcal{B} = \{b_\ell, b_m, b_h\}$ with
$b_\ell < b_m < b_h$, and define the per-question expected reward:
\begin{equation}
  R(q, b) \;=\; P\!\left(\textrm{correct} \mid q, b\right) - \lambda \cdot
  \frac{b}{b_h},
  \label{eq:reward}
\end{equation}
where $\lambda \ge 0$ penalizes token cost.
Suppose (i)~$P(\textrm{correct} \mid q, b)$ is non-decreasing in $b$
with diminishing returns (concavity), and
(ii)~the marginal accuracy gain of increasing the budget is
non-increasing in $p(q)$ (easy questions benefit less from extra budget).
Then the optimal policy $\pi^*(q) = \argmax_{b \in \mathcal{B}} R(q, b)$
takes the form of a \emph{threshold policy on $C_K$}:
\begin{equation}
  \pi^*(q) =
  \begin{cases}
    b_\ell & \text{if } C_K(q) \ge \tau_1, \\
    b_m    & \text{if } \tau_2 \le C_K(q) < \tau_1, \\
    b_h    & \text{if } C_K(q) < \tau_2,
  \end{cases}
  \label{eq:threshold-policy}
\end{equation}
for thresholds $0 \le \tau_2 \le \tau_1 \le 1$ that depend on $\lambda$
and the accuracy--budget curve.
\end{proposition}

\begin{proof}[Proof sketch]
Under assumptions~(i) and (ii), the marginal value of upgrading from
$b_\ell$ to $b_m$ (or $b_m$ to $b_h$) is a non-increasing function of
$p(q)$.
By Proposition~\ref{prop:monotone}, $C_K$ is a monotone proxy for
$p(q)$, so the indifference points where the marginal accuracy gain
equals the marginal cost $\lambda(b_{k+1}-b_k)/b_h$ correspond to
thresholds on $C_K$.
Between these thresholds, the optimal action is constant.
Concavity of the accuracy--budget curve ensures that at most three
regimes arise, matching the three routes in \method{}.
A formal derivation via Lagrangian relaxation is in
Appendix~\ref{app:proofs}.
\end{proof}

\paragraph{Connection to \method{} design.}
Proposition~\ref{prop:threshold} provides a normative justification for
the three-route architecture of \method{}:
high consensus ($C_K \ge \tau_1$) triggers the \textbf{Accept} route
(easy),
partial consensus ($\tau_2 \le C_K < \tau_1$) triggers the \textbf{Extend}
route (medium),
and low consensus ($C_K < \tau_2$) triggers the \textbf{Deliberate} route
(hard).
The thresholds $\tau_1, \tau_2$ are the only hyperparameters introduced by
\method{}, and they can be set by a lightweight sweep on a small
validation set or derived from the cost parameter $\lambda$.
Crucially, this result shows that no finer partition of the consensus
space is needed under the stated regularity conditions---the three-route
design is \emph{optimal} within the class of consensus-based policies.


We now provide a theoretical foundation for the central design choice in
\method{}: using multi-path consensus rather than single-path features as
the difficulty signal for budget allocation.
Let $q$ denote a question, $D(q)\in\{\textrm{easy},\textrm{medium},\textrm{hard}\}$
its latent difficulty, $X_i$ the answer produced by the $i$-th independent
reasoning path, and
$C_K \triangleq \frac{1}{\binom{K}{2}}\sum_{i<j}\mathbf{1}[X_i = X_j]$
the pairwise agreement ratio among $K$ paths.

%% ===== Proposition 1 =====
\begin{proposition}[Consensus is more informative than a single path]
\label{prop:info}
For any $K \ge 2$,
\begin{equation}
  I\!\left(D;\, C_K\right) \;\ge\; I\!\left(D;\, X_1\right),
  \label{eq:mi-bound}
\end{equation}
where $I(\cdot\,;\cdot)$ denotes mutual information.
The inequality is strict whenever there exist two difficulty levels
$d, d'$ such that the answer distributions $P(X_1\!\mid\! D\!=\!d)$
and $P(X_1\!\mid\! D\!=\!d')$ differ.
\end{proposition}

\begin{proof}[Proof sketch]
By assumption, paths are conditionally i.i.d.\ given $q$ and hence given
$D(q)$.
Consider the joint sample $\mathbf{X}_K = (X_1, \dots, X_K)$.
By the data-processing inequality applied to the Markov chain
$D \to \mathbf{X}_K \to C_K$, we have
$I(D; C_K) \le I(D; \mathbf{X}_K)$.
In the reverse direction, $X_1$ is a deterministic function of
$\mathbf{X}_K$, so
$I(D; X_1) \le I(D; \mathbf{X}_K)$.
The key step is to observe that $C_K$ is a symmetric function of all $K$
samples and therefore captures \emph{inter-path} variation that a single
sample cannot.
Formally, $C_K$ is a function of the empirical distribution
$\hat{P}_K = \frac{1}{K}\sum_i \delta_{X_i}$, which is a sufficient
statistic for the family
$\{P(\mathbf{X}_K \!\mid\! D\!=\!d)\}_{d}$ by the factorization theorem
(since $P(\mathbf{X}_K \!\mid\! D) = \prod_i P(X_i \!\mid\! D)$
depends on $\mathbf{X}_K$ only through $\hat{P}_K$).
Sufficiency implies $I(D; \hat{P}_K) = I(D; \mathbf{X}_K)$,
and because $C_K$ is an injective function of $\hat{P}_K$ when the answer
space has at most $K$ distinct values, we obtain
$I(D; C_K) = I(D; \mathbf{X}_K) \ge I(D; X_1)$.
Strictness follows because $K \ge 2$ independent samples from distinct
distributions yield strictly more Fisher information and hence strictly
more mutual information than a single sample, provided
$P(X_1 \!\mid\! D\!=\!d) \ne P(X_1 \!\mid\! D\!=\!d')$ for some
$d \ne d'$.
A complete proof is given in Appendix~\ref{app:proofs}.
\end{proof}

%% ===== Proposition 2 =====
\begin{proposition}[Consensus as a monotone difficulty estimator]
\label{prop:monotone}
Let $p(q) = P(\text{correct}\!\mid\! q)$ denote the per-question
accuracy of a single path.
If paths are conditionally independent given $q$ and the answer space
is binary (correct/incorrect), then:
\begin{enumerate}[nosep,leftmargin=*]
  \item $\mathbb{E}[C_K \mid q]$ is strictly increasing in $p(q)$ for
        $p(q) > \tfrac{1}{2}$.
  \item As $K \to \infty$,
    \begin{equation}
      C_K \;\xrightarrow{\;a.s.\;}\; p(q)^{2} + \bigl(1 - p(q)\bigr)^{2},
      \label{eq:consensus-limit}
    \end{equation}
    which is a strictly decreasing function of difficulty
    (lower $p(q) \Rightarrow$ lower $C_\infty$) on $p(q)\in(\tfrac12,1]$.
  \item For finite $K$, $\mathrm{Var}(C_K \mid q) = O(1/K)$, so the
        difficulty estimate concentrates rapidly.
\end{enumerate}
Consequently, if $P(\textrm{correct}\!\mid\!\textrm{easy}) >
P(\textrm{correct}\!\mid\!\textrm{hard})$, then
$\mathbb{E}[C_K \mid \textrm{easy}] > \mathbb{E}[C_K \mid \textrm{hard}]$.
\end{proposition}

\begin{proof}[Proof sketch]
Under conditional independence and a binary answer space,
$C_K = \binom{K}{2}^{-1}\sum_{i<j}\mathbf{1}[X_i = X_j]$
has expectation
$\mathbb{E}[C_K \mid q] = p(q)^2 + (1-p(q))^2$.
Differentiating with respect to $p$:
$\frac{d}{dp}[p^2 + (1-p)^2] = 2(2p - 1) > 0$ for $p > \tfrac{1}{2}$.
Statement~(2) follows from the strong law of large numbers applied to
$\hat{P}_K$, and~(3) from the bounded-differences inequality applied
to $C_K$ (changing one path shifts $C_K$ by at most $2/(K-1)$).
\end{proof}

\begin{remark}
When the answer space is multiclass (e.g., numerical answers with many
possible values), the correct-answer probability $p(q)$ still dominates
$C_K$ because incorrect answers are typically dispersed across many
values.
Our experiments confirm this: on GSM8K, the average number of distinct
incorrect answers per question is $3.2$, ensuring that consensus is
driven primarily by agreement on the correct answer.
\end{remark}

%% ===== Proposition 3 =====
\begin{proposition}[Fundamental limitation of single-path signals]
\label{prop:single-path}
In the black-box setting (no access to logits, hidden states, or
internal model representations), the only observable from a single
reasoning path is the generated text $X_1$.
For any difficulty estimator $f\colon \mathcal{X} \to \mathbb{R}$
that is a deterministic function of $X_1$:
\begin{equation}
  I\!\left(D;\, f(X_1)\right) \;\le\; I\!\left(D;\, X_1\right)
  \;\le\; I\!\left(D;\, C_K\right) \quad \text{for } K \ge 2.
  \label{eq:single-path-bound}
\end{equation}
\end{proposition}

\begin{proof}[Proof sketch]
The first inequality is the data-processing inequality applied to
$D \to X_1 \to f(X_1)$.
The second follows from Proposition~\ref{prop:info}.
The bound is tight only in the degenerate case where a single sample
fully determines $D$---which requires the answer distributions under
different difficulty levels to have disjoint support, a condition that
fails in practice (easy and hard questions can produce the same answer
text).
\end{proof}

\paragraph{Empirical validation.}
We compute the Spearman rank correlation between candidate difficulty
signals and the oracle difficulty label (fraction of budgets on which
the model errs).
Single-path features---answer presence, token utilization, and reasoning
coherence---achieve $|\rho| \in [0.08, 0.19]$.
In contrast, the consensus signal $C_K$ with $K = 8$ achieves
$\rho = -0.51$ ($p < 10^{-88}$).
Multi-path consensus also yields a 72.8\% precision in identifying
easy-vs-hard questions, closing 73\% of the gap between single-path
heuristics and the oracle classifier.
These numbers corroborate Propositions~\ref{prop:info}--\ref{prop:single-path}:
consensus is not merely incrementally better; it provides a qualitatively
different signal.

%% ===== Proposition 4: Budget allocation =====
\begin{proposition}[Optimal budget allocation is a threshold policy on consensus]
\label{prop:threshold}
Consider a budget set $\mathcal{B} = \{b_\ell, b_m, b_h\}$ with
$b_\ell < b_m < b_h$, and define the per-question expected reward:
\begin{equation}
  R(q, b) \;=\; P\!\left(\textrm{correct} \mid q, b\right) - \lambda \cdot
  \frac{b}{b_h},
  \label{eq:reward}
\end{equation}
where $\lambda \ge 0$ penalizes token cost.
Suppose (i)~$P(\textrm{correct} \mid q, b)$ is non-decreasing in $b$
with diminishing returns (concavity), and
(ii)~the marginal accuracy gain of increasing the budget is
non-increasing in $p(q)$ (easy questions benefit less from extra budget).
Then the optimal policy $\pi^*(q) = \argmax_{b \in \mathcal{B}} R(q, b)$
takes the form of a \emph{threshold policy on $C_K$}:
\begin{equation}
  \pi^*(q) =
  \begin{cases}
    b_\ell & \text{if } C_K(q) \ge \tau_1, \\
    b_m    & \text{if } \tau_2 \le C_K(q) < \tau_1, \\
    b_h    & \text{if } C_K(q) < \tau_2,
  \end{cases}
  \label{eq:threshold-policy}
\end{equation}
for thresholds $0 \le \tau_2 \le \tau_1 \le 1$ that depend on $\lambda$
and the accuracy--budget curve.
\end{proposition}

\begin{proof}[Proof sketch]
Under assumptions~(i) and (ii), the marginal value of upgrading from
$b_\ell$ to $b_m$ (or $b_m$ to $b_h$) is a non-increasing function of
$p(q)$.
By Proposition~\ref{prop:monotone}, $C_K$ is a monotone proxy for
$p(q)$, so the indifference points where the marginal accuracy gain
equals the marginal cost $\lambda(b_{k+1}-b_k)/b_h$ correspond to
thresholds on $C_K$.
Between these thresholds, the optimal action is constant.
Concavity of the accuracy--budget curve ensures that at most three
regimes arise, matching the three routes in \method{}.
A formal derivation via Lagrangian relaxation is in
Appendix~\ref{app:proofs}.
\end{proof}

\paragraph{Connection to \method{} design.}
Proposition~\ref{prop:threshold} provides a normative justification for
the three-route architecture of \method{}:
high consensus ($C_K \ge \tau_1$) triggers the \textbf{Accept} route
(easy),
partial consensus ($\tau_2 \le C_K < \tau_1$) triggers the \textbf{Extend}
route (medium),
and low consensus ($C_K < \tau_2$) triggers the \textbf{Deliberate} route
(hard).
The thresholds $\tau_1, \tau_2$ are the only hyperparameters introduced by
\method{}, and they can be set by a lightweight sweep on a small
validation set or derived from the cost parameter $\lambda$.
Crucially, this result shows that no finer partition of the consensus
space is needed under the stated regularity conditions---the three-route
design is \emph{optimal} within the class of consensus-based policies.


We now provide a theoretical foundation for the central design choice in
\method{}: using multi-path consensus rather than single-path features as
the difficulty signal for budget allocation.
Let $q$ denote a question, $D(q)\in\{\textrm{easy},\textrm{medium},\textrm{hard}\}$
its latent difficulty, $X_i$ the answer produced by the $i$-th independent
reasoning path, and
$C_K \triangleq \frac{1}{\binom{K}{2}}\sum_{i<j}\mathbf{1}[X_i = X_j]$
the pairwise agreement ratio among $K$ paths.

%% ===== Proposition 1 =====
\begin{proposition}[Consensus is more informative than a single path]
\label{prop:info}
For any $K \ge 2$,
\begin{equation}
  I\!\left(D;\, C_K\right) \;\ge\; I\!\left(D;\, X_1\right),
  \label{eq:mi-bound}
\end{equation}
where $I(\cdot\,;\cdot)$ denotes mutual information.
The inequality is strict whenever there exist two difficulty levels
$d, d'$ such that the answer distributions $P(X_1\!\mid\! D\!=\!d)$
and $P(X_1\!\mid\! D\!=\!d')$ differ.
\end{proposition}

\begin{proof}[Proof sketch]
By assumption, paths are conditionally i.i.d.\ given $q$ and hence given
$D(q)$.
Consider the joint sample $\mathbf{X}_K = (X_1, \dots, X_K)$.
By the data-processing inequality applied to the Markov chain
$D \to \mathbf{X}_K \to C_K$, we have
$I(D; C_K) \le I(D; \mathbf{X}_K)$.
In the reverse direction, $X_1$ is a deterministic function of
$\mathbf{X}_K$, so
$I(D; X_1) \le I(D; \mathbf{X}_K)$.
The key step is to observe that $C_K$ is a symmetric function of all $K$
samples and therefore captures \emph{inter-path} variation that a single
sample cannot.
Formally, $C_K$ is a function of the empirical distribution
$\hat{P}_K = \frac{1}{K}\sum_i \delta_{X_i}$, which is a sufficient
statistic for the family
$\{P(\mathbf{X}_K \!\mid\! D\!=\!d)\}_{d}$ by the factorization theorem
(since $P(\mathbf{X}_K \!\mid\! D) = \prod_i P(X_i \!\mid\! D)$
depends on $\mathbf{X}_K$ only through $\hat{P}_K$).
Sufficiency implies $I(D; \hat{P}_K) = I(D; \mathbf{X}_K)$,
and because $C_K$ is an injective function of $\hat{P}_K$ when the answer
space has at most $K$ distinct values, we obtain
$I(D; C_K) = I(D; \mathbf{X}_K) \ge I(D; X_1)$.
Strictness follows because $K \ge 2$ independent samples from distinct
distributions yield strictly more Fisher information and hence strictly
more mutual information than a single sample, provided
$P(X_1 \!\mid\! D\!=\!d) \ne P(X_1 \!\mid\! D\!=\!d')$ for some
$d \ne d'$.
A complete proof is given in Appendix~\ref{app:proofs}.
\end{proof}

%% ===== Proposition 2 =====
\begin{proposition}[Consensus as a monotone difficulty estimator]
\label{prop:monotone}
Let $p(q) = P(\text{correct}\!\mid\! q)$ denote the per-question
accuracy of a single path.
If paths are conditionally independent given $q$ and the answer space
is binary (correct/incorrect), then:
\begin{enumerate}[nosep,leftmargin=*]
  \item $\mathbb{E}[C_K \mid q]$ is strictly increasing in $p(q)$ for
        $p(q) > \tfrac{1}{2}$.
  \item As $K \to \infty$,
    \begin{equation}
      C_K \;\xrightarrow{\;a.s.\;}\; p(q)^{2} + \bigl(1 - p(q)\bigr)^{2},
      \label{eq:consensus-limit}
    \end{equation}
    which is a strictly decreasing function of difficulty
    (lower $p(q) \Rightarrow$ lower $C_\infty$) on $p(q)\in(\tfrac12,1]$.
  \item For finite $K$, $\mathrm{Var}(C_K \mid q) = O(1/K)$, so the
        difficulty estimate concentrates rapidly.
\end{enumerate}
Consequently, if $P(\textrm{correct}\!\mid\!\textrm{easy}) >
P(\textrm{correct}\!\mid\!\textrm{hard})$, then
$\mathbb{E}[C_K \mid \textrm{easy}] > \mathbb{E}[C_K \mid \textrm{hard}]$.
\end{proposition}

\begin{proof}[Proof sketch]
Under conditional independence and a binary answer space,
$C_K = \binom{K}{2}^{-1}\sum_{i<j}\mathbf{1}[X_i = X_j]$
has expectation
$\mathbb{E}[C_K \mid q] = p(q)^2 + (1-p(q))^2$.
Differentiating with respect to $p$:
$\frac{d}{dp}[p^2 + (1-p)^2] = 2(2p - 1) > 0$ for $p > \tfrac{1}{2}$.
Statement~(2) follows from the strong law of large numbers applied to
$\hat{P}_K$, and~(3) from the bounded-differences inequality applied
to $C_K$ (changing one path shifts $C_K$ by at most $2/(K-1)$).
\end{proof}

\begin{remark}
When the answer space is multiclass (e.g., numerical answers with many
possible values), the correct-answer probability $p(q)$ still dominates
$C_K$ because incorrect answers are typically dispersed across many
values.
Our experiments confirm this: on GSM8K, the average number of distinct
incorrect answers per question is $3.2$, ensuring that consensus is
driven primarily by agreement on the correct answer.
\end{remark}

%% ===== Proposition 3 =====
\begin{proposition}[Fundamental limitation of single-path signals]
\label{prop:single-path}
In the black-box setting (no access to logits, hidden states, or
internal model representations), the only observable from a single
reasoning path is the generated text $X_1$.
For any difficulty estimator $f\colon \mathcal{X} \to \mathbb{R}$
that is a deterministic function of $X_1$:
\begin{equation}
  I\!\left(D;\, f(X_1)\right) \;\le\; I\!\left(D;\, X_1\right)
  \;\le\; I\!\left(D;\, C_K\right) \quad \text{for } K \ge 2.
  \label{eq:single-path-bound}
\end{equation}
\end{proposition}

\begin{proof}[Proof sketch]
The first inequality is the data-processing inequality applied to
$D \to X_1 \to f(X_1)$.
The second follows from Proposition~\ref{prop:info}.
The bound is tight only in the degenerate case where a single sample
fully determines $D$---which requires the answer distributions under
different difficulty levels to have disjoint support, a condition that
fails in practice (easy and hard questions can produce the same answer
text).
\end{proof}

\paragraph{Empirical validation.}
We compute the Spearman rank correlation between candidate difficulty
signals and the oracle difficulty label (fraction of budgets on which
the model errs).
Single-path features---answer presence, token utilization, and reasoning
coherence---achieve $|\rho| \in [0.08, 0.19]$.
In contrast, the consensus signal $C_K$ with $K = 8$ achieves
$\rho = -0.51$ ($p < 10^{-88}$).
Multi-path consensus also yields a 72.8\% precision in identifying
easy-vs-hard questions, closing 73\% of the gap between single-path
heuristics and the oracle classifier.
These numbers corroborate Propositions~\ref{prop:info}--\ref{prop:single-path}:
consensus is not merely incrementally better; it provides a qualitatively
different signal.

%% ===== Proposition 4: Budget allocation =====
\begin{proposition}[Optimal budget allocation is a threshold policy on consensus]
\label{prop:threshold}
Consider a budget set $\mathcal{B} = \{b_\ell, b_m, b_h\}$ with
$b_\ell < b_m < b_h$, and define the per-question expected reward:
\begin{equation}
  R(q, b) \;=\; P\!\left(\textrm{correct} \mid q, b\right) - \lambda \cdot
  \frac{b}{b_h},
  \label{eq:reward}
\end{equation}
where $\lambda \ge 0$ penalizes token cost.
Suppose (i)~$P(\textrm{correct} \mid q, b)$ is non-decreasing in $b$
with diminishing returns (concavity), and
(ii)~the marginal accuracy gain of increasing the budget is
non-increasing in $p(q)$ (easy questions benefit less from extra budget).
Then the optimal policy $\pi^*(q) = \argmax_{b \in \mathcal{B}} R(q, b)$
takes the form of a \emph{threshold policy on $C_K$}:
\begin{equation}
  \pi^*(q) =
  \begin{cases}
    b_\ell & \text{if } C_K(q) \ge \tau_1, \\
    b_m    & \text{if } \tau_2 \le C_K(q) < \tau_1, \\
    b_h    & \text{if } C_K(q) < \tau_2,
  \end{cases}
  \label{eq:threshold-policy}
\end{equation}
for thresholds $0 \le \tau_2 \le \tau_1 \le 1$ that depend on $\lambda$
and the accuracy--budget curve.
\end{proposition}

\begin{proof}[Proof sketch]
Under assumptions~(i) and (ii), the marginal value of upgrading from
$b_\ell$ to $b_m$ (or $b_m$ to $b_h$) is a non-increasing function of
$p(q)$.
By Proposition~\ref{prop:monotone}, $C_K$ is a monotone proxy for
$p(q)$, so the indifference points where the marginal accuracy gain
equals the marginal cost $\lambda(b_{k+1}-b_k)/b_h$ correspond to
thresholds on $C_K$.
Between these thresholds, the optimal action is constant.
Concavity of the accuracy--budget curve ensures that at most three
regimes arise, matching the three routes in \method{}.
A formal derivation via Lagrangian relaxation is in
Appendix~\ref{app:proofs}.
\end{proof}

\paragraph{Connection to \method{} design.}
Proposition~\ref{prop:threshold} provides a normative justification for
the three-route architecture of \method{}:
high consensus ($C_K \ge \tau_1$) triggers the \textbf{Accept} route
(easy),
partial consensus ($\tau_2 \le C_K < \tau_1$) triggers the \textbf{Extend}
route (medium),
and low consensus ($C_K < \tau_2$) triggers the \textbf{Deliberate} route
(hard).
The thresholds $\tau_1, \tau_2$ are the only hyperparameters introduced by
\method{}, and they can be set by a lightweight sweep on a small
validation set or derived from the cost parameter $\lambda$.
Crucially, this result shows that no finer partition of the consensus
space is needed under the stated regularity conditions---the three-route
design is \emph{optimal} within the class of consensus-based policies.
